\documentclass[11pt]{article}
\usepackage{amssymb,amsmath,amsthm,multicol,enumitem,graphicx,mathtools}
\setlength{\oddsidemargin}{0in}
\setlength{\textwidth}{6.5in}
\setlength{\topmargin}{-.35in}
\setlength{\textheight}{9in}
\usepackage{xcolor}

\parindent 0pt
\parskip 7pt

\DeclareMathOperator{\im}{Im}
\DeclareMathOperator{\preim}{PreIm}
\DeclareMathOperator{\lcm}{lcm}
\DeclareMathOperator{\id}{id}
\DeclareMathOperator{\ord}{ord}
\DeclarePairedDelimiter\ceil{\lceil}{\rceil}
\DeclarePairedDelimiter\floor{\lfloor}{\rfloor}

\setlist[itemize,1]{label=$\bullet$}
\setlist[itemize,2]{label={\tiny $\blacktriangleright$}}
\setlist[itemize,3]{label=$\diamond$}
\setlist[itemize,4]{label=$\circ$}

\setlist[enumerate]{itemsep=15pt}
\setlist[itemize]{itemsep=5pt}
\addtolength{\jot}{5pt}

\usepackage{titlesec}
\titlespacing*{\subsection}{0pt}{25pt}{0pt}

\newif\ifSHOWSOLUTION
\SHOWSOLUTIONtrue

\newcommand{\showsolution}[1]{\par
  \ifSHOWSOLUTION
  \textbf{\textcolor{blue}{Solution.}} {\color{blue}#1}
  \fi
}

\begin{document}
\begin{center}
  {\LARGE 21-127 Set Theory

    \large Yuxiang Huang [yh4]

  September 20th, 2026}
\end{center}

Prove that for any sets $A$ and $B$,
$$
A \subseteq B \text{ or } B \subseteq A
\quad \Longleftrightarrow \quad
\mathcal{P}(A) \cup \mathcal{P}(B) = \mathcal{P}(A \cup B)
$$

\newpage
\showsolution{
  Let $A$ and $B$ be sets.
  \begin{itemize}
    \item[$(\Rightarrow):$] Assume $A \subseteq B$ or
      $B \subseteq A$.
      \begin{itemize}
        \item[$(\subseteq):$] Let
          $X \in \mathcal{P}(A) \cup \mathcal{P}(B)$.
          By def of set union, we have
          $X \in \mathcal{P}(A)$ or $X \in \mathcal{P}(B)$.
          By def of power set, we have
          $X \subseteq A$ or $X \subseteq B$.
          In either case, we have
          $X \subseteq A \cup B$. Therefore,
          $X \in \mathcal{P}(A \cup B)$.

        \item[$(\supseteq):$] Let $X \in \mathcal{P}(A \cup B)$.
          By def of power set, we have
          $X \subseteq A \cup B$.
          \begin{itemize}
            \item Case 1: $A \subseteq B$. Then
              $A \cup B = B$, so $X \subseteq B$.
              By def of power set, we have
              $X \in \mathcal{P}(B)$. By def of set union, we have
              $X \in \mathcal{P}(A) \cup \mathcal{P}(B)$.

            \item Case 2: $B \subseteq A$. Then
              $A \cup B = A$, so $X \subseteq A$. By def of power set, we have
              $X \in \mathcal{P}(A)$. By def of set union, we have
              $X \in \mathcal{P}(A) \cup \mathcal{P}(B)$.
          \end{itemize}
          Thus, in either case, $X \in \mathcal{P}(A) \cup \mathcal{P}(B)$.
      \end{itemize}

    \item[$(\Leftarrow):$] Assume
      $\mathcal{P}(A) \cup \mathcal{P}(B) = \mathcal{P}(A \cup B)$.

      AFSOC $A \nsubseteq B$ and
      $B \nsubseteq A$. Then there exist $a \in A \setminus B$ and
      $b \in B \setminus A$. Since $\{a,b\} \subseteq A \cup B$,
      by def of power set, we have
      $\{a,b\} \in \mathcal{P}(A \cup B)$.

      However, since
      $b \in B \setminus A$, by def of set difference, we have
      $b \notin A$. By def of subset, we have
      $\{a,b\} \nsubseteq A$. By def of power set, we have
      $\{a,b\} \notin \mathcal{P}(A)$.
      Similarly, since
      $a \in A \setminus B$, by def of set difference, we have
      $a \notin B$. By def of subset, we have
      $\{a,b\} \nsubseteq B$. By def of power set, we have
      $\{a,b\} \notin \mathcal{P}(B)$.
      Therefore,
      since $\{a,b\} \notin \mathcal{P}(A)$ and
      $\{a,b\} \notin \mathcal{P}(B)$,
      by def of set union, we have that
      $\{a,b\} \notin \mathcal{P}(A) \cup \mathcal{P}(B)$.

      This contradicts the assumption that
      $\mathcal{P}(A) \cup \mathcal{P}(B) = \mathcal{P}(A \cup B)$.
      Therefore, we conclude that $A \subseteq B$ or $B \subseteq A$.
      $\Rightarrow\Leftarrow$
  \end{itemize}
}

\end{document}
